\documentclass[letterpaper,11pt]{article}

\usepackage{latexsym}
\usepackage[empty]{fullpage}
\usepackage{titlesec}
\usepackage{marvosym}
\usepackage[usenames,dvipsnames]{color}
\usepackage{verbatim}
\usepackage{enumitem}
\usepackage[hidelinks]{hyperref}
\usepackage{fancyhdr}
\usepackage[english]{babel}
\usepackage[utf8]{inputenc}
\usepackage{tabularx}
\input{glyphtounicode}

\pagestyle{fancy}
\fancyhf{}
\fancyfoot{}
\renewcommand{\headrulewidth}{0pt}
\renewcommand{\footrulewidth}{0pt}

\addtolength{\oddsidemargin}{-0.5in}
\addtolength{\evensidemargin}{-0.5in}
\addtolength{\textwidth}{1in}
\addtolength{\topmargin}{-0.88in}
\addtolength{\textheight}{1.9in}

\urlstyle{same}
\raggedbottom
\raggedright
\setlength{\tabcolsep}{0in}

\titleformat{\section}{
  \vspace{-4pt}\scshape\raggedright\large
}{}{0em}{}[\color{black}\titlerule \vspace{-5pt}]

\pdfgentounicode=1
\usepackage{textgreek}
\newcommand{\resumeName}[1]{\textbf{\Huge \scshape #1} \\ \vspace{1pt}}

\newcommand{\resumeItem}[1]{
  \item\small{
    {#1 \vspace{-2pt}}
  }
}

\newcommand{\resumeSubheading}[4]{
  \vspace{-2pt}\item
    \begin{tabular*}{0.97\textwidth}[t]{l@{\extracolsep{\fill}}r}
      \textbf{#1} & #2 \\
      \textit{\small#3} & \textit{\small #4} \\
    \end{tabular*}\vspace{-7pt}
}

\newcommand{\resumeSubSubheading}[2]{
    \item
    \begin{tabular*}{0.97\textwidth}{l@{\extracolsep{\fill}}r}
      \textit{\small#1} & \textit{\small #2} \\
    \end{tabular*}\vspace{-7pt}
}

\newcommand{\resumeProjectHeading}[2]{
    \item
    \begin{tabular*}{0.97\textwidth}{l@{\extracolsep{\fill}}r}
      \small#1 & #2 \\
    \end{tabular*}\vspace{-7pt}
}

\newcommand{\resumePatentHeading}[3]{
    \item
    \begin{tabularx}{0.97\textwidth}[t]{@{}>{\raggedright\arraybackslash}X r@{}}
      \small\textbf{#1} -- \textit{#2} & \small #3 \\
    \end{tabularx}\vspace{-7pt}
}

\newcommand{\resumeSubItem}[1]{\resumeItem{#1}\vspace{-4pt}}

\renewcommand\labelitemii{$\vcenter{\hbox{\tiny$\bullet$}}$}

\newcommand{\resumeSubHeadingListStart}{\begin{itemize}[leftmargin=0.15in, label={}]}
\newcommand{\resumeSubHeadingListEnd}{\end{itemize}}
\newcommand{\resumeItemListStart}{\begin{itemize}[itemsep=-1pt,topsep=2pt]}
\newcommand{\resumeItemListEnd}{\end{itemize}\vspace{-5pt}}

\begin{document}

\begin{center}
  \resumeName{Harinandan Kotamsetti}
    \small Cupertino, CA $|$ +1 (669) 264-9458 $|$ \href{mailto:hareee234@gmail.com}{\underline{hareee234@gmail.com}} $|$
    \href{https://www.linkedin.com/in/hareee234/}{\underline{linkedin.com/in/hareee234}} $|$
    \href{https://github.com/AwesomenessReborn}{\underline{github.com/AwesomenessReborn}}
\end{center}

\section{Education}
  \resumeSubHeadingListStart
    \resumeSubheading
      {\textbf{San Jose State University}}{San Jose, CA}
      {Bachelor of Science in Computer Engineering}{Expected Dec 2026}
      \resumeItemListStart
        \resumeItem{\textbf{Coursework:} OS, Computer Architecture, C/C++, FPGA, Verilog, Linear Algebra, x86 \& MIPS Assembly}
      \resumeItemListEnd
  \resumeSubHeadingListEnd

\section{Experience}
  \resumeSubHeadingListStart
    \resumeSubheading
      {\textbf{STMicroelectronics}}{Santa Clara, CA}
      {\textbf{Software \& Algorithm Intern}}{Mar 2025 -- Present}
      \resumeItemListStart
        \resumeItem{
          \textbf{Built the on-device Android inference path for a walking-speed model} (129 KB CNN-LSTM, 0.059 m/s MAE): ring-buffered C++ IMU preprocessing at 208 Hz (replacing the Python pipeline) and a JNI bridge to Qualcomm's NPU runtime, running in real time at $\sim$15 ms per 2.5 s window; \textbf{demoed at CES 2026}.
        }
        \resumeItem{
          \textbf{Ported an 887K-parameter VR head-motion model to Qualcomm's Hexagon NPU} (PyTorch $\rightarrow$ ONNX $\rightarrow$ QAIRT on HTP v81), matching PyTorch output within 0.056\textdegree{} on-device after fixing converter and tensor-layout failures.
        }
        \resumeItem{
          Deployed always-on fall detection and wrist-tilt recognition entirely on ST's ISPU/MLC with \textbf{zero host-MCU compute}, sustaining 104 Hz sensor fusion at 3.08 ms worst case in $\sim$20 KB code / 1.1 KB RAM.
        }
      \resumeItemListEnd

    \resumeSubheading
      {\textbf{Spartan Racing (Formula SAE)}}{San Jose, CA}
      {\textbf{Embedded Systems \& Software Developer}}{Aug 2023 -- Feb 2025}
      \resumeItemListStart
        \resumeItem{
          \textbf{Built a Python hardware-in-the-loop test suite} on a breakout board emulating an accumulator module, validating OV/UV thresholds and fault logic without a live 400 V pack.
        }
        \resumeItem{
          \textbf{Drove component and trace changes across 4 BMS board revisions}, reviewing schematics with the EE team and feeding back design issues found during firmware bring-up.
        }
        \resumeItem{
          \textbf{Cut STM32F105 BMS polling latency 14$\times$ (7.0 s $\rightarrow$ 0.5 s)} across a 768-cell, 400 V pack by removing blocking CAN/thermistor delays, replacing 10 ms transmit waits with 3-mailbox availability checks, and moving to static buffers.
        }
        \resumeItem{
          \textbf{Implemented interrupt-driven BMS fault handling with TIM7/NVIC ISRs}, encoding fault/warning states into 8-bit status bytes for deterministic HV shutdown signaling across the 8-module accumulator.
        }
        \resumeItem{
          \textbf{Migrated STM32 BMS firmware from LTC6811 to LTC6813} across a \textbf{768-cell pack}, updating ADC/register handling, static buffers, and CAN encoding for 16-bit voltage/temperature telemetry.
        }
      \resumeItemListEnd
  \resumeSubHeadingListEnd

\section{Patents \& Invention Disclosures}
  \resumeSubHeadingListStart
    \resumePatentHeading
      {Motion-to-Photon Latency Compensation for XR/AR/VR Glasses}
      {STMicroelectronics, Co-inventor}
      {Invention Disclosure, 2026}
    \resumePatentHeading
      {Glasses with Wellness and Health Monitoring}
      {STMicroelectronics, Co-inventor}
      {Filed May 2026}
  \resumeSubHeadingListEnd

\section{Projects}
  \resumeSubHeadingListStart
    \resumeProjectHeading
      {\textbf{MERIT UAV --- Autonomous Surveillance Drone} $|$ \emph{Python, pymavlink, YOLOv8, CUDA, Jetson}}{2026}
      \resumeItemListStart
        \resumeItem{
          \textbf{Integrated and flight-tested a Jetson Orin Nano with a Pixhawk over 921600-baud UART}, commanding takeoff, altitude, position, and yaw via MAVLink in GUIDED mode.
        }
        \resumeItem{
          Ran a \textbf{concurrent 4 Hz GPS telemetry daemon alongside 30 FPS CUDA YOLOv8n detection}, geotagging detections and gating uploads on GPS lock; vision pipeline ground-tested.
        }
      \resumeItemListEnd

    \resumeProjectHeading
      {\textbf{MIPS SoC with Hardware Accelerator} $|$ \emph{Verilog, Xilinx Vivado, Basys 3 FPGA}}{2024}
      \resumeItemListStart
        \resumeItem{
          Designed and deployed \textbf{32-bit MIPS SoC} on Artix-7 FPGA with custom single-cycle CPU, memory-mapped GPIO, and hardware factorial accelerator using Moore FSM with datapath/control separation.
        }
      \resumeItemListEnd
  \resumeSubHeadingListEnd

\section{Technical Skills}
  \resumeSubHeadingListStart
    \small{\item{
      \textbf{Embedded \& Hardware}{: STM32 (ARM Cortex-M3), Device Drivers, SPI, isoSPI, I2C, CAN, BLE, MEMS IMUs, Real-Time Systems} \\
      \textbf{Testing \& Debugging}{: HIL Testing, Regression Testing, Performance Benchmarking, JTAG Debuggers, Oscilloscopes, Software Simulation} \\
      \textbf{Systems \& Deployment}{: C, C++, Python, JNI/NDK, Android, Linux, CMake, Make, Git, STM32CubeIDE} \\
      \textbf{EDA \& Signal Processing}{: Verilog, Altium, LTspice, MATLAB, Simulink, Xilinx Vivado, DSP, Signal Processing} \\
      \textbf{ML Deployment}{: PyTorch, ONNX, QNN/QAIRT, TFLite, TensorFlow}
    }}
  \resumeSubHeadingListEnd

\end{document}
